\section{Alignement final DevSecOps}

\subsection{Position de reference}

Le scenario officiel de soutenance de \textit{SecureRAG Hub} est le mode
\texttt{demo}. Jenkins reste l'autorite CI/CD officielle du projet, tandis que
GitHub Actions est conserve comme historique technique. Cette decision evite de
presenter deux chaines concurrentes et clarifie le perimetre de demonstration.

\subsection{Elements validates}

Les elements suivants sont consideres comme validates dans l'etat courant du
projet :

\begin{itemize}
  \item execution des tests applicatifs de demonstration ;
  \item couverture de tests ;
  \item linting Python avec Ruff ;
  \item analyse Semgrep ;
  \item detection de secrets avec Gitleaks ;
  \item scan filesystem avec Trivy ;
  \item deploiement local sur \texttt{kind} ;
  \item overlay Kubernetes \texttt{demo} ;
  \item Jenkins local Docker ;
  \item campagne CD \texttt{demo} en mode \texttt{dry-run} ;
  \item generation d'un support pack de soutenance.
\end{itemize}

\subsection{Elements prets mais dependants de l'environnement}

Les composants suivants sont implementes et scriptes, mais leur validation
complete depend de l'environnement d'execution :

\begin{itemize}
  \item declenchement Jenkins par webhook GitHub ;
  \item generation de SBOM par \texttt{Syft} ;
  \item signature et verification d'images par \texttt{Cosign} ;
  \item promotion par digest sans reconstruction ;
  \item deploiement verifie en mode \texttt{execute} ;
  \item activation de \texttt{metrics-server} ;
  \item observation effective des \textit{HorizontalPodAutoscaler} ;
  \item installation de \texttt{Kyverno} ;
  \item passage progressif des policies de \texttt{Audit} vers \texttt{Enforce}.
\end{itemize}

Ces points ne doivent pas etre presentes comme executes si les artefacts de
preuve correspondants ne sont pas disponibles. Ils doivent etre presentes comme
prets a etre rejoues, avec dependances explicites.

\subsection{Preuves a conserver}

Les preuves finales attendues sont les suivantes :

\begin{itemize}
  \item \texttt{artifacts/jenkins/github-webhook-validation.md} ;
  \item \texttt{artifacts/release/supply-chain-execute-summary.md} ;
  \item \texttt{artifacts/release/sign-summary.txt} ;
  \item \texttt{artifacts/release/verify-summary.txt} ;
  \item \texttt{artifacts/release/promotion-by-digest-summary.txt} ;
  \item \texttt{artifacts/release/promotion-digests.txt} ;
  \item \texttt{artifacts/sbom/sbom-index.txt} ;
  \item \texttt{artifacts/validation/cluster-security-addons.md} ;
  \item \texttt{artifacts/final/final-validation-summary.md} ;
  \item \texttt{artifacts/support-pack/*.tar.gz}.
\end{itemize}

\subsection{Formulation recommandee pour la soutenance}

La formulation recommandee est la suivante :

\begin{quote}
Le socle DevSecOps, Kubernetes et demonstration \texttt{demo} est operationnel
et valide. La supply chain avancee est structuree et executable, mais sa preuve
complete depend de la disponibilite locale de Docker, du registre OCI, des cles
Cosign, de Syft et des images signees. Les addons de securite cluster
\texttt{metrics-server} et \texttt{Kyverno} sont installables et observables via
des scripts dedies, avec un passage prudent de \texttt{Audit} vers
\texttt{Enforce}.
\end{quote}

